\documentclass[12pt, a4paper]{book}

% ─── Core packages ────────────────────────────────────────────────────────────
\usepackage[T1]{fontenc}
\usepackage[utf8]{inputenc}
\usepackage[charter, cal=cmcal]{mathdesign}
\usepackage[scaled=0.92]{helvet}
% microtype removed
\usepackage{geometry}
\geometry{
  a4paper,
  top=30mm, bottom=30mm,
  inner=32mm, outer=25mm,
  marginparwidth=18mm
}

% ─── Mathematics ──────────────────────────────────────────────────────────────
\usepackage{amsmath, amsthm}

% ─── Chemistry packages ───────────────────────────────────────────────────────
\usepackage[version=4]{mhchem}          % \ce{H2O}, reaction arrows, etc.
% chemfig removed                    % structural formulas
% bohr removed                       % Bohr atomic diagrams
% elements removed                   % element data

% ─── Graphics & colour ────────────────────────────────────────────────────────
\usepackage{graphicx}
\usepackage[dvipsnames,svgnames,table]{xcolor}
\usepackage{tikz}
\usetikzlibrary{arrows.meta, decorations.pathmorphing, shapes.geometric,
                positioning, calc, backgrounds, shadows}
\usepackage{pgfplots}
\pgfplotsset{compat=1.18}
\usepackage{tcolorbox}
\tcbuselibrary{skins, breakable, theorems}

% ─── Layout & typography ──────────────────────────────────────────────────────
\usepackage{titlesec}
\usepackage{titletoc}
\usepackage{fancyhdr}
\usepackage{enumitem}
\usepackage{booktabs}
\usepackage{array}
\usepackage{multicol}
\usepackage{marginnote}
\usepackage{mdframed}
\usepackage{wrapfig}
\usepackage{caption}
\usepackage{subcaption}
\usepackage{setspace}
\usepackage{parskip}
\usepackage{siunitx}

% ─── Hyperlinks ───────────────────────────────────────────────────────────────
\usepackage[colorlinks=true, linkcolor=BrickRed,
            urlcolor=NavyBlue, citecolor=ForestGreen]{hyperref}

% ══════════════════════════════════════════════════════════════════════════════
% COLOUR PALETTE
% ══════════════════════════════════════════════════════════════════════════════
\definecolor{JourneyColor}{HTML}{065F46}   % forest green — journey headings
\definecolor{ModuleColor}{HTML}{0F766E}    % deep teal — module headings
\definecolor{PuzzleColor}{HTML}{B45309}    % deep amber — puzzle boxes
\definecolor{DiscoveryColor}{HTML}{D97706} % golden amber — discovery boxes
\definecolor{WarningColor}{HTML}{E11D48}   % coral red — "model breaks" boxes
\definecolor{ProblemColor}{HTML}{0F766E}   % dark teal — problem sets
\definecolor{ReflectColor}{HTML}{047857}   % green — reflection boxes
\definecolor{HistoryColor}{HTML}{9A3412}   % brown — history notes
\definecolor{LightBg}{HTML}{F8FAFC}        % near-white background        % near-white background

% ══════════════════════════════════════════════════════════════════════════════
% TCOLORBOX ENVIRONMENTS
% ══════════════════════════════════════════════════════════════════════════════

%% Puzzle box ─────────────────────────────────────────────────────────────────
\newtcolorbox{puzzlebox}[1][]{
  enhanced, breakable,
  colback=PuzzleColor!6, colframe=PuzzleColor,
  fonttitle=\bfseries\sffamily,
  title={The Puzzle},
  attach boxed title to top left={yshift=-2mm, xshift=6mm},
  boxed title style={colback=PuzzleColor, colframe=PuzzleColor},
  borderline west={3pt}{0pt}{PuzzleColor},
  #1
}

%% Discovery box ──────────────────────────────────────────────────────────────
\newtcolorbox{discoverybox}[1][]{
  enhanced, breakable,
  colback=DiscoveryColor!8, colframe=DiscoveryColor,
  fonttitle=\bfseries\sffamily,
  title={Discovery},
  attach boxed title to top left={yshift=-2mm, xshift=6mm},
  boxed title style={colback=DiscoveryColor, colframe=DiscoveryColor},
  borderline west={3pt}{0pt}{DiscoveryColor},
  #1
}

%% Problem set box ────────────────────────────────────────────────────────────
\newtcolorbox{problembox}[1][]{
  enhanced, breakable,
  colback=ProblemColor!5, colframe=ProblemColor,
  fonttitle=\bfseries\sffamily,
  title={Problems},
  attach boxed title to top left={yshift=-2mm, xshift=6mm},
  boxed title style={colback=ProblemColor, colframe=ProblemColor},
  borderline west={3pt}{0pt}{ProblemColor},
  #1
}

%% Deep problems box ──────────────────────────────────────────────────────────
\newtcolorbox{deepbox}[1][]{
  enhanced, breakable,
  colback=JourneyColor!5, colframe=JourneyColor,
  fonttitle=\bfseries\sffamily,
  title={Deep Problems},
  attach boxed title to top left={yshift=-2mm, xshift=6mm},
  boxed title style={colback=JourneyColor, colframe=JourneyColor},
  borderline west={4pt}{0pt}{JourneyColor},
  #1
}

%% History note ───────────────────────────────────────────────────────────────
\newtcolorbox{historybox}[1][]{
  enhanced, breakable,
  colback=HistoryColor!7, colframe=HistoryColor!60,
  fonttitle=\bfseries\sffamily\small,
  title={Historical Note},
  attach boxed title to top left={yshift=-2mm, xshift=6mm},
  boxed title style={colback=HistoryColor!80, colframe=HistoryColor},
  #1
}

%% Reflection box ─────────────────────────────────────────────────────────────
\newtcolorbox{reflectionbox}[1][]{
  enhanced, breakable,
  colback=ReflectColor!6, colframe=ReflectColor,
  fonttitle=\bfseries\sffamily,
  title={Reflection},
  attach boxed title to top left={yshift=-2mm, xshift=6mm},
  boxed title style={colback=ReflectColor, colframe=ReflectColor},
  borderline west={3pt}{0pt}{ReflectColor},
  #1
}

%% Model-breaks warning ───────────────────────────────────────────────────────
\newtcolorbox{modelbreaks}[1][]{
  enhanced, breakable,
  colback=WarningColor!6, colframe=WarningColor,
  fonttitle=\bfseries\sffamily,
  title={Where This Model Breaks},
  attach boxed title to top left={yshift=-2mm, xshift=6mm},
  boxed title style={colback=WarningColor, colframe=WarningColor},
  borderline west={4pt}{0pt}{WarningColor},
  #1
}

%% Named concept box ──────────────────────────────────────────────────────────
\newtcolorbox{namedconcept}[2][]{
  enhanced,
  colback=ModuleColor!6, colframe=ModuleColor,
  fonttitle=\bfseries\sffamily\large,
  title={#2},
  attach boxed title to top left={yshift=-2mm, xshift=6mm},
  boxed title style={colback=ModuleColor, colframe=ModuleColor},
  #1
}

% ══════════════════════════════════════════════════════════════════════════════
% THEOREM-LIKE ENVIRONMENTS
% ══════════════════════════════════════════════════════════════════════════════
\theoremstyle{definition}
\newtheorem{observation}{Observation}[chapter]
\newtheorem{question}{Question}[chapter]
\newtheorem{experiment}{Experiment}[chapter]
\newtheorem{pattern}{Pattern}[chapter]

% ══════════════════════════════════════════════════════════════════════════════
% SECTION FORMATTING
% ══════════════════════════════════════════════════════════════════════════════
\titleformat{\chapter}[display]
  {\normalfont\sffamily\huge\bfseries\color{JourneyColor}}
  {\textcolor{JourneyColor!60}{\chaptertitlename\ \thechapter}}
  {10pt}{\Huge}[\vspace{2pt}\textcolor{JourneyColor}{\titlerule[1.5pt]}]

\titleformat{\section}
  {\normalfont\sffamily\Large\bfseries\color{ModuleColor}}
  {\thesection}{1em}{}[\textcolor{ModuleColor!40}{\titlerule[0.6pt]}]

\titleformat{\subsection}
  {\normalfont\sffamily\large\bfseries\color{JourneyColor!80}}
  {\thesubsection}{1em}{}

\titleformat{\subsubsection}
  {\normalfont\sffamily\normalsize\bfseries}
  {\thesubsubsection}{1em}{}

% ══════════════════════════════════════════════════════════════════════════════
% HEADER / FOOTER
% ══════════════════════════════════════════════════════════════════════════════
\pagestyle{fancy}
\fancyhf{}
\fancyhead[LE]{\small\sffamily\textcolor{JourneyColor}{\leftmark}}
\fancyhead[RO]{\small\sffamily\textcolor{ModuleColor}{\rightmark}}
\fancyfoot[C]{\small\thepage}
\renewcommand{\headrulewidth}{0.4pt}

% ══════════════════════════════════════════════════════════════════════════════
% CHEMFIG SETTINGS
% ══════════════════════════════════════════════════════════════════════════════

% ══════════════════════════════════════════════════════════════════════════════
% CUSTOM COMMANDS
% ══════════════════════════════════════════════════════════════════════════════
\newcommand{\journeyname}[1]{\textcolor{JourneyColor}{\textbf{\sffamily#1}}}
\newcommand{\chem}[1]{\ce{#1}}
\newcommand{\important}[1]{\textcolor{ModuleColor}{\textbf{#1}}}
\newcommand{\newterm}[1]{\textit{\textcolor{DiscoveryColor}{#1}}}

% margin annotation
\newcommand{\sidenote}[1]{%
  \marginnote{\footnotesize\sffamily\textcolor{JourneyColor!70}{#1}}}

% ══════════════════════════════════════════════════════════════════════════════
% DOCUMENT BEGINS
% ══════════════════════════════════════════════════════════════════════════════
\begin{document}

% ─── Title page ──────────────────────────────────────────────────────────────
\begin{titlepage}
\pagecolor{JourneyColor}
\color{white}
\centering
\vspace*{3cm}

{\fontsize{11}{13}\selectfont\sffamily\MakeUppercase{A Discovery Curriculum}}

\vspace{0.6cm}
{\fontsize{42}{46}\selectfont\bfseries\sffamily The Chemistry\\[6pt] of Everything}

\vspace{0.8cm}
\textcolor{white!60}{\rule{0.5\textwidth}{1pt}}

\vspace{0.8cm}
{\large\sffamily From Substances to Structure to the World They Build}

\vspace{3cm}

\begin{tcolorbox}[
  enhanced, width=0.75\textwidth,
  colback=white!10, colframe=white!40,
  boxrule=0.8pt, arc=4pt
]
\centering\sffamily\small
This document is written for the student who will not merely\\
\textit{use} chemistry as a tool, but who will one day\\
\textit{extend} it. Every concept you meet here was once unknown.\\
You are invited to rediscover it.
\end{tcolorbox}

\vfill
{\sffamily\small Kenya Mathematical Olympiad Programme \quad\textbullet\quad Discovery Series}
\pagecolor{white}
\end{titlepage}

% ─── TOC ─────────────────────────────────────────────────────────────────────
\tableofcontents
\clearpage

% ══════════════════════════════════════════════════════════════════════════════
%  JOURNEY 1 — WHAT IS STUFF?
% ══════════════════════════════════════════════════════════════════════════════
\chapter[Journey 1 — What Is Stuff?]
        {Journey One\\[8pt]\textcolor{ModuleColor}{\large What Is Stuff?\\ Matter and Its Patterns}}
\label{ch:journey1}

\begin{tcolorbox}[
  enhanced, colback=JourneyColor!8, colframe=JourneyColor,
  arc=4pt, left=6mm, right=6mm, top=4mm, bottom=4mm
]
\sffamily\small
\textbf{What this journey is about.}
Chemistry begins not with atoms — atoms are already an answer.
It begins with a much older question:
\textit{why do different substances behave so differently?}
Why does wood burn while stone does not?
Why does salt dissolve in water but oil refuses?
Why does heating copper ore in a fire produce
a shining metal while the ore itself is dull and brittle?

In this journey we do what the earliest chemists did:
we observe, we classify, we measure, and we chase
the simplest possible explanation for stubborn facts.
By the end, we will have \textit{invented} the concept of the atom ---
not because someone told us atoms exist,
but because the data leave us no other choice.
\end{tcolorbox}

\vspace{1em}

\begin{center}
\begin{tikzpicture}[
  node distance=1.6cm,
  every node/.style={font=\sffamily\small},
  arrow/.style={-{Stealth[length=5pt]}, thick, color=JourneyColor!70}
]
  \node[draw=JourneyColor, fill=JourneyColor!10, rounded corners=4pt,
        minimum width=3.2cm, minimum height=0.8cm, align=center]
        (obs) {Observation:\\substances differ};
  \node[draw=ModuleColor, fill=ModuleColor!10, rounded corners=4pt,
        minimum width=3.2cm, minimum height=0.8cm, align=center,
        right=of obs]
        (clas) {Classification:\\elements \& compounds};
  \node[draw=DiscoveryColor, fill=DiscoveryColor!10, rounded corners=4pt,
        minimum width=3.2cm, minimum height=0.8cm, align=center,
        right=of clas]
        (cons) {Conservation:\\mass is conserved};
  \node[draw=PuzzleColor, fill=PuzzleColor!10, rounded corners=4pt,
        minimum width=3.2cm, minimum height=0.8cm, align=center,
        right=of cons]
        (atom) {Fixed ratios:\\the atomic hypothesis};

  \draw[arrow] (obs)  -- (clas);
  \draw[arrow] (clas) -- (cons);
  \draw[arrow] (cons) -- (atom);
\end{tikzpicture}
\end{center}

\vspace{1.5em}

% ══════════════════════════════════════════════════════════════════════════════
%  MODULE 1.1 — THE STUBBORN QUESTION
% ══════════════════════════════════════════════════════════════════════════════
\section{Module 1.1 — The Stubborn Question: Why Do Substances Differ?}
\label{sec:mod1.1}

\subsection*{Motivation}

Before we discuss any chemistry at all, we must sit with a fact that is
easy to overlook because we are so accustomed to it:

\begin{center}
\begin{tcolorbox}[
  enhanced, colback=LightBg, colframe=JourneyColor!40,
  width=0.88\textwidth, arc=3pt
]
\centering\sffamily
The world is made of an enormous variety of substances.\\[4pt]
Iron is hard and conducts electricity.
Sulfur is brittle and does not.\\
Water flows and extinguishes fire.
Ethanol flows and feeds fire.\\
Salt is a white crystal that your body needs to survive.\\
Arsenic is a white powder that will kill you.\\[4pt]
\textit{Same colour. Very different substance.}
\end{tcolorbox}
\end{center}

\vspace{0.5em}

Why should matter come in so many distinct varieties?
This is the question chemistry was born to answer.
And the answer, when it comes, will be surprisingly deep:
the properties of a substance are consequences of its \emph{composition}
and \emph{structure} at the very smallest scale.

But we are not there yet.
First, we must learn to look carefully.

% ─── Section: Observation ────────────────────────────────────────────────────
\subsection{The Art of Careful Observation}

The earliest move in chemistry is not calculation.
It is \newterm{systematic observation} — noting not just that
substances differ, but \emph{how} and \emph{when} they differ.

Consider the following observations, which any careful person
in the ancient world could have made.

\begin{observation}[Fire and transformation]
When wood burns, the solid wood disappears and is replaced by
a small amount of ash, heat, and smoke.
The product is nothing like the original wood.
\label{obs:burning}
\end{observation}

\begin{observation}[Heating ore]
When greenish \emph{malachite} (a copper-containing rock)
is heated strongly with charcoal,
a shining red-orange metal appears.
The rock becomes metal.
\label{obs:ore}
\end{observation}

\begin{observation}[Dissolving]
When common salt is placed in water, it disappears.
The water now tastes salty.
When the water is evaporated, the salt reappears —
identical to what was added.
\label{obs:dissolving}
\end{observation}

\begin{observation}[Fixed composition]
No matter where in the world you find pure water,
it always has the same properties: same boiling point,
same density, same taste, same behaviour.
Water from the Nile and water from the Amazon are \emph{indistinguishable}.
\label{obs:fixedcomp}
\end{observation}

\vspace{0.5em}

These four observations are profoundly different in character,
and it took centuries to understand why.
Observations~\ref{obs:burning} and~\ref{obs:ore} involve
\emph{transformation} — the substance changes permanently.
Observation~\ref{obs:dissolving} involves \emph{mixing} —
the salt returns unchanged.
Observation~\ref{obs:fixedcomp} is perhaps the most striking:
pure substances have \emph{fixed, reproducible properties}.

\begin{puzzlebox}
You have a clear liquid in a flask.
You do not know what it is.

\begin{enumerate}[leftmargin=*, label=(\alph*)]
  \item List five \emph{physical} measurements you could make to
        characterise it, without chemically altering it.
  \item What does it mean for a substance to be \emph{pure}?
        How would you test for purity using only the properties
        you listed?
  \item Two liquids are both clear and colourless.
        One is water, one is ethanol.
        They look identical.
        Design an experiment (without tasting) to distinguish them.
\end{enumerate}
\end{puzzlebox}

% ─── Section: Physical vs Chemical ───────────────────────────────────────────
\subsection{Physical Properties vs Chemical Behaviour}

The earliest classification in chemistry distinguishes
two types of properties:

\begin{itemize}[leftmargin=2em]
  \item \newterm{Physical properties} can be measured without
        changing what the substance fundamentally is:
        colour, mass, density, boiling point, melting point,
        electrical conductivity.
  \item \newterm{Chemical properties} describe how a substance
        interacts with other substances — what it reacts with,
        what it becomes, what energy is released or absorbed.
\end{itemize}

\vspace{0.5em}

\begin{center}
\begin{tikzpicture}[
  every node/.style={font=\sffamily\small},
  arrow/.style={-{Stealth[length=4pt]}, thick}
]
  % Two columns: Physical | Chemical
  \node[draw=ModuleColor, fill=ModuleColor!8, rounded corners=4pt,
        minimum width=5cm, minimum height=1cm, align=center]
        (phys) at (0,0) {\textbf{Physical Property}\\
                          substance unchanged after measurement};
  \node[draw=WarningColor, fill=WarningColor!8, rounded corners=4pt,
        minimum width=5cm, minimum height=1cm, align=center]
        (chem) at (7.5,0) {\textbf{Chemical Property}\\
                             substance is transformed};

  % Examples below
  \node[align=left, text=ModuleColor] at (0,-1.8) {
    density, boiling point,\\
    colour, conductivity,\\
    hardness, solubility};
  \node[align=left, text=WarningColor] at (7.5,-1.8) {
    flammability, reactivity\\
    with acid, ability to\\
    rust, toxicity};
\end{tikzpicture}
\end{center}

\vspace{0.5em}

Why does this distinction matter?
Because it tells us something deep:
\emph{a substance can be characterised by its properties,
and no two distinct pure substances have identical properties}.
Properties are like fingerprints.
This will become the foundation of chemical identification.

\begin{problembox}
\begin{enumerate}[label=\textbf{\arabic*.}]
  \item Classify each property as physical or chemical:
    \begin{multicols}{2}
    \begin{itemize}[nosep]
      \item Melting point of ice: \SI{0}{\celsius}
      \item Iron rusts in moist air
      \item Gold density: \SI{19.3}{\gram\per\centi\metre\cubed}
      \item Magnesium burns with bright white flame
      \item Ethanol boils at \SI{78}{\celsius}
      \item Sodium reacts violently with water
    \end{itemize}
    \end{multicols}

  \item A student claims that dissolving salt in water is a
        \emph{chemical} change because the salt disappears.
        Another claims it is \emph{physical} because the salt
        can be recovered by evaporation.
        Who is correct, and why?

  \item Consider ice, liquid water, and steam.
        Do they have the same physical properties?
        The same chemical properties?
        What does your answer tell you about the
        relationship between physical state and chemical identity?
\end{enumerate}
\end{problembox}

% ─── Section: Mixtures vs Pure Substances ────────────────────────────────────
\subsection{Mixtures and Pure Substances}

A key early discovery is that not all matter is equally \emph{uniform}.

\begin{observation}[Heterogeneous vs homogeneous matter]
Gravel mixed with sand looks non-uniform to the eye.
Salt dissolved in water looks completely uniform.
Yet both are \emph{mixtures} that can be separated by physical means.
\end{observation}

\vspace{0.5em}

This leads to our first classification tree:

\begin{center}
\begin{tikzpicture}[
  every node/.style={font=\sffamily\small, align=center},
  edge/.style={thick, color=JourneyColor!70},
  level 1/.style={sibling distance=6cm},
  level 2/.style={sibling distance=3cm}
]
  \node[draw=JourneyColor, fill=JourneyColor!10,
        rounded corners=4pt, minimum width=3cm, minimum height=0.7cm]
        {Matter}
    child[edge from parent/.style={draw=JourneyColor, thick}]{
      node[draw=ModuleColor, fill=ModuleColor!10,
           rounded corners=3pt, minimum width=2.8cm]{Pure\\Substance}
      child{ node[draw=ModuleColor!70, fill=ModuleColor!5,
                  rounded corners=2pt]{Element} }
      child{ node[draw=ModuleColor!70, fill=ModuleColor!5,
                  rounded corners=2pt]{Compound} }
    }
    child[edge from parent/.style={draw=JourneyColor, thick}]{
      node[draw=DiscoveryColor!80, fill=DiscoveryColor!10,
           rounded corners=3pt, minimum width=2.8cm]{Mixture}
      child{ node[draw=DiscoveryColor!60, fill=DiscoveryColor!5,
                  rounded corners=2pt]{Homogeneous\\(solution)} }
      child{ node[draw=DiscoveryColor!60, fill=DiscoveryColor!5,
                  rounded corners=2pt]{Heterogeneous} }
    };
\end{tikzpicture}
\end{center}

\vspace{0.5em}

The most important boundary here is between a \newterm{pure substance}
and a \newterm{mixture}.

\begin{namedconcept}{Pure Substance}
A pure substance has \emph{fixed, definite composition}
and \emph{fixed properties} that do not depend on its origin or
how it was prepared.

A pure substance cannot be separated into simpler components
by \emph{physical} means (filtering, evaporation, distillation).
\end{namedconcept}

\vspace{0.5em}

Notice what we did not say: we said nothing about atoms.
This definition is \emph{operational} — it is defined by what you
can and cannot do to the substance.
This is exactly how early chemists thought.

\begin{puzzlebox}
Consider the following and classify each:

\begin{center}
\begin{tabular}{lll}
\toprule
Substance & Observation & Your classification \\
\midrule
Seawater & Tastes salty; evaporation leaves solid & ? \\
Distilled water & Always boils at exactly \SI{100}{\celsius} & ? \\
Bronze & Gold-coloured metal, variable hardness & ? \\
Dry air & Uniform gas; \ce{N2} extractable by cooling & ? \\
Diamond & Always same density, hardness & ? \\
Crude oil & Distillation gives many different liquids & ? \\
\bottomrule
\end{tabular}
\end{center}

After classifying, answer:
\begin{enumerate}[label=(\alph*)]
  \item What test distinguishes a pure substance from a mixture?
  \item Can a mixture be \emph{uniform} in appearance?
        Does uniformity guarantee purity?
\end{enumerate}
\end{puzzlebox}

% ─── Section: Elements ───────────────────────────────────────────────────────
\subsection{Elements: Substances That Cannot Be Broken Down Further}

We can divide pure substances further.
Some can be decomposed into simpler substances by chemical means.
Others cannot.

\begin{namedconcept}{Element}
An \newterm{element} is a pure substance that \emph{cannot be broken
down into simpler substances} by any chemical process.

Historically, something was declared an element when no known
chemical procedure could decompose it further.
\end{namedconcept}

\vspace{0.5em}

\begin{namedconcept}{Compound}
A \newterm{compound} is a pure substance that \emph{can be decomposed}
into simpler substances, and whose properties are entirely
\emph{different} from the elements that compose it.
\end{namedconcept}

\vspace{0.5em}

The distinction sounds simple. It was extraordinarily hard to establish in practice. 

Consider water.
For over two thousand years, water was believed to be an \emph{element}.
It took until 1781 for Henry Cavendish to show that burning hydrogen
gas in air produces water — demonstrating that water is a \emph{compound}.

\begin{historybox}
\textbf{Cavendish and the Composition of Water (1781)}

Henry Cavendish was a reclusive English scientist who avoided
company but was meticulous with measurements.
He discovered that a gas he called \emph{inflammable air}
(what we now call hydrogen) burns to produce water.
He later showed the ratio: water is approximately
\SI{11.1}{\percent} hydrogen and \SI{88.9}{\percent} oxygen by mass.

This single experiment overturned millennia of belief.
Water was not an element.
It was built from two things that were.

\begin{center}
\ce{2H2 + O2 -> 2H2O}
\end{center}
\end{historybox}

% ─── Section: The Strangeness of Compounds ───────────────────────────────────
\subsection{The Strangeness of Compounds}

Here is something that should bother you.

Sodium is a \emph{silvery metal} that reacts explosively with water.
Chlorine is a \emph{yellow-green poisonous gas}.

When they react together:

\vspace{0.3em}
\begin{center}
\ce{2Na_{(s)} + Cl2_{(g)} -> 2NaCl_{(s)}}
\end{center}
\vspace{0.3em}

The product is common table salt — a white crystal essential for life,
that dissolves harmlessly in water.

\begin{center}
\begin{tikzpicture}[every node/.style={font=\sffamily\small}]
  % Sodium
  \node[draw=JourneyColor, fill=JourneyColor!10, rounded corners=3pt,
        minimum width=2.5cm, minimum height=1.8cm, align=center]
        (na) at (0,0) {\textbf{Sodium}\\silver metal\\reacts with water};
  % Plus sign
  \node at (3,0) {\Large +};
  % Chlorine
  \node[draw=WarningColor, fill=WarningColor!10, rounded corners=3pt,
        minimum width=2.5cm, minimum height=1.8cm, align=center]
        (cl) at (5,0) {\textbf{Chlorine}\\toxic green gas\\corrosive};
  % Arrow
  \draw[-{Stealth[length=6pt]}, thick, color=ModuleColor] (6.5,0) -- (8.5,0)
        node[midway, above, font=\sffamily\small]{react};
  % Salt
  \node[draw=ModuleColor, fill=ModuleColor!10, rounded corners=3pt,
        minimum width=2.8cm, minimum height=1.8cm, align=center]
        (salt) at (10.5,0) {\textbf{Sodium chloride}\\white crystal\\edible salt};
\end{tikzpicture}
\end{center}

\vspace{0.5em}

This is one of the most important lessons in chemistry:
\emph{the properties of a compound are not the sum of the properties
of its elements}.
When elements combine to form a compound,
something genuinely new is created.

A compound is not a mixture of its elements.
It is a new substance entirely.

\begin{deepbox}
\begin{enumerate}[label=\textbf{D\arabic*.}]
  \item Iron (\ce{Fe}) is a strong metal. Sulfur (\ce{S}) is a
        yellow brittle powder. When heated together they form
        iron(II) sulfide (\ce{FeS}), a grey-black brittle solid.
        \begin{enumerate}[label=(\alph*)]
          \item A magnet attracts iron but not \ce{FeS}.
                Is this consistent with \ce{FeS} being a compound?
                Explain.
          \item If you mixed iron filings and sulfur powder
                \emph{without} heating, would a magnet separate them?
                Why does this distinguish a mixture from a compound?
        \end{enumerate}

  \item Water (\ce{H2O}) and hydrogen peroxide (\ce{H2O2})
        have the same elements: hydrogen and oxygen.
        Water is stable and life-sustaining.
        Hydrogen peroxide is a strong oxidising agent that can
        bleach hair and disinfect wounds.
        \begin{enumerate}[label=(\alph*)]
          \item What does this tell you about the relationship
                between composition and properties?
          \item Is the \emph{ratio} of elements in a compound
                fixed or variable? How would you test this experimentally?
        \end{enumerate}

  \item A student decomposes \SI{18.0}{\gram} of water by passing
        electricity through it (electrolysis).
        She collects \SI{2.0}{\gram} of hydrogen gas and
        \SI{16.0}{\gram} of oxygen gas.
        She repeats with \SI{36.0}{\gram} of water.
        \begin{enumerate}[label=(\alph*)]
          \item Predict the masses of hydrogen and oxygen produced.
          \item What pattern do you see?
          \item Another student claims the ratio might change if
                the water comes from a different source.
                Design an experiment to test this claim.
        \end{enumerate}
\end{enumerate}
\end{deepbox}

% ─── Section: Separation Methods ─────────────────────────────────────────────
\subsection{Separating Mixtures: Physical Methods}

An important experimental skill — one that shaped the history
of chemistry — is the ability to separate mixtures into their components.

The key principle: \emph{different components of a mixture differ in
their physical properties}, and we can exploit these differences
to separate them without any chemical reaction occurring.

\begin{center}
\begin{tikzpicture}[
  every node/.style={font=\sffamily\footnotesize, align=center},
  arrow/.style={-{Stealth[length=4pt]}, thick, color=JourneyColor!60}
]
  % Central node
  \node[draw=JourneyColor, fill=JourneyColor!10, rounded corners=5pt,
        minimum width=3cm, minimum height=1cm, font=\sffamily\small\bfseries]
        (mix) {Separating\\Mixtures};

  % Surrounding methods
  \node[draw=ModuleColor!80, fill=ModuleColor!8, rounded corners=3pt,
        minimum width=2.8cm, minimum height=0.9cm]
        (filt) at (0, 2.8) {Filtration\\solid from liquid};
  \node[draw=ModuleColor!80, fill=ModuleColor!8, rounded corners=3pt,
        minimum width=2.8cm, minimum height=0.9cm]
        (evap) at (3.5, 1.8) {Evaporation\\dissolved solid from solution};
  \node[draw=ModuleColor!80, fill=ModuleColor!8, rounded corners=3pt,
        minimum width=2.8cm, minimum height=0.9cm]
        (dist) at (3.5, -1.8) {Distillation\\liquids by boiling point};
  \node[draw=ModuleColor!80, fill=ModuleColor!8, rounded corners=3pt,
        minimum width=2.8cm, minimum height=0.9cm]
        (chrom) at (0, -2.8) {Chromatography\\by solubility};
  \node[draw=ModuleColor!80, fill=ModuleColor!8, rounded corners=3pt,
        minimum width=2.8cm, minimum height=0.9cm]
        (mag) at (-3.5, -1.8) {Magnetism\\magnetic from non-magnetic};
  \node[draw=ModuleColor!80, fill=ModuleColor!8, rounded corners=3pt,
        minimum width=2.8cm, minimum height=0.9cm]
        (crys) at (-3.5, 1.8) {Crystallisation\\by solubility difference};

  \foreach \n in {filt, evap, dist, chrom, mag, crys}
    \draw[arrow] (mix) -- (\n);
\end{tikzpicture}
\end{center}

\vspace{0.5em}

Each separation method relies on a \emph{physical property difference}:
filtration uses particle size;
distillation uses boiling point;
chromatography uses relative solubility.
None of these methods create new substances.

\begin{problembox}
\begin{enumerate}[label=\textbf{\arabic*.}]
  \item A mixture contains salt (\ce{NaCl}), sand (\ce{SiO2}),
        and iron filings.
        \begin{enumerate}[label=(\alph*)]
          \item Design a three-step procedure to obtain
                each pure component.
          \item State which physical property difference
                each step exploits.
        \end{enumerate}

  \item Crude oil is a mixture of many hydrocarbons.
        They can be separated by heating and collecting vapours
        at different temperatures.
        \begin{enumerate}[label=(\alph*)]
          \item What physical property is this separation based on?
          \item If two components have very similar properties,
                how would you modify the procedure?
        \end{enumerate}

  \item Ink is a mixture of coloured compounds.
        If you place a dot of ink on filter paper and allow
        water to seep through, the colours separate.
        \begin{enumerate}[label=(\alph*)]
          \item What does this tell you about the ink's composition?
          \item Could this technique tell you whether a substance
                is pure or mixed? Explain.
        \end{enumerate}
\end{enumerate}
\end{problembox}

% ─── Section: Towards the Atom ───────────────────────────────────────────────
\subsection{The Question at the Edge of Observation}

We have established:
substances differ in properties;
some are pure, some are mixtures;
some pure substances can be decomposed (compounds);
others cannot (elements).

This leaves us with a question that observation alone cannot answer:

\begin{center}
\begin{tcolorbox}[
  enhanced, colback=PuzzleColor!8, colframe=PuzzleColor,
  width=0.82\textwidth, arc=3pt,
  fontupper=\sffamily\large\bfseries, halign=center
]
What is an element \emph{made of}?\\[4pt]
\normalfont\normalsize\sffamily
If we could divide a piece of iron into smaller and smaller pieces,\\
would we eventually reach something that is not iron?\\
Or would every piece, however small, still be iron?
\end{tcolorbox}
\end{center}

\vspace{0.5em}

This question is ancient.
The Greek philosopher Democritus (around 400 BCE)
proposed that matter is made of \emph{indivisible particles}
he called \emph{atomos} — uncuttable.
But this was pure philosophy, with no experimental evidence.

The answer would have to wait two thousand years,
for a set of precise measurements that pointed inescapably
toward a hidden structure.

That is the subject of Module 1.4.
Before we get there, we need one more observation
from the history of chemistry — one that transformed
the entire field.

% ─── Section: Reflection ─────────────────────────────────────────────────────
\subsection*{Reflection}

\begin{reflectionbox}
Take ten minutes to write your answers to the following.
You can revise your answers after completing the later modules.

\begin{enumerate}[label=\textbf{R\arabic*.}]
  \item We built a classification tree: matter splits into
        mixtures and pure substances; pure substances split into
        elements and compounds.
        \textbf{What observations forced each branch of this classification?}
        In other words, what experiments would you do to decide
        where an unknown substance belongs?

  \item The properties of sodium chloride are completely different
        from those of sodium metal and chlorine gas.
        \textbf{What does this tell us about the nature of chemical bonding?}
        (You do not need to answer this fully now — just write
        your current intuition.)

  \item We defined an element as a substance that cannot be decomposed
        by chemical means.
        \textbf{What is the weakness of this definition?}
        (Hint: what would happen if we hadn't yet discovered
        all possible chemical reactions?)

  \item A pure substance always has fixed properties,
        regardless of its origin.
        \textbf{Why is this fact so important?}
        What would chemistry look like if the same substance
        had different properties depending on where it was found?
\end{enumerate}
\end{reflectionbox}

\vspace{1.5em}

\begin{center}
\begin{tcolorbox}[
  enhanced, colback=DiscoveryColor!8, colframe=DiscoveryColor,
  width=0.88\textwidth, arc=3pt
]
\centering\sffamily\small
\textbf{What we invented in this module.}\\[4pt]
We invented the idea that substances can be \emph{classified by
what can and cannot be done to them}.
This classification — into mixtures, compounds, and elements —
was not handed to us by nature. It was a \emph{choice of framework},
and it turned out to be the right one.
Chemists did not discover elements by finding tiny objects called elements.
They discovered them by finding substances that \emph{resisted}
all attempts at decomposition.\\[4pt]
When we later discover that elements correspond to distinct types of atoms,
we will see that this operational definition was pointing at
something deeply real — but we did not need to know that yet.
\end{tcolorbox}
\end{center}

\vspace{2em}
\noindent\textcolor{JourneyColor!40}{\rule{\textwidth}{0.6pt}}
\begin{center}
\textit{\sffamily\small
Module 1.1 complete. Continue to Module 1.2:
\textbf{Burning, Dissolving, Reacting — The First Clues}.}
\end{center}

\end{document}